% main.tex - Documento raiz de la propuesta competitiva / PPC (memoria tecnica).
% La estructura del documento la DICTA LA CONVOCATORIA: ajusta secciones,
% orden y limites de paginas a las bases de la convocatoria (rag/calls/) y
% verifica la cobertura en context/convocatoria_checklist.md. Compilar con
% pdflatex + bibtex.
\documentclass[12pt]{article}

\usepackage[utf8]{inputenc}
\usepackage[T1]{fontenc}
\usepackage[margin=2.5cm]{geometry}
\usepackage{booktabs}
\usepackage{hyperref}

% Marcadores academicos \todo{...} (ver .roo/rules/00_academic_core.md):
% fallback compilable. Si anades el paquete todonotes, elimina esta linea.
\providecommand{\todo}[1]{\textbf{[TODO: #1]}}

\title{TODO: titulo de la propuesta (proyecto de colaboracion publico-privada)}
\author{TODO: coordinador/a del consorcio}
\date{TODO: fecha de envio}

\begin{document}

\maketitle

% Resumen ejecutivo (seccion sin numerar): la estructura de 9 secciones es
% orientativa; la convocatoria manda (context/convocatoria_checklist.md).
\input{src/00_resumen}
\input{src/01_contexto}
\input{src/02_objetivos}
\input{src/03_excellence}
\input{src/04_impact}
\input{src/05_wps}
\input{src/06_consortium}
\input{src/07_budget}
\input{src/08_risks}

\bibliographystyle{plain}
\bibliography{src/references}

\end{document}
